\daytitle{30}{综合训练与后续学习路线}

\begin{learningbox}{本节导读}
最后一节不再新增大量概念，而是把前面分散学到的知识转化为可执行的方法。综合题不是把 30 节公式堆在一起，而是先读题、画图、选对象、分过程，再为每一段选择合适规律并复核结果。本节用力学综合、电磁综合和实验综合收束全书，同时给出完成一个月入门后的继续学习路线。
\end{learningbox}

\textbf{本节读题路线：} 综合题第一遍先补完整图示和情境信息；第二遍只找状态变化、区域变化和相互作用变化，把全过程切成若干阶段。列式时每个方程旁边标出规律名称和适用阶段，最后用量纲、极限和方向检查结果是否可信。

\section*{一、学习目标}
\begin{enumerate}
  \item 能按“五步拆解法”读综合题：补图、选对象、分过程、列规律、复核结果。
  \item 能按条件改变的位置，把复杂过程拆成若干基本模型。
  \item 会在方程旁写出规律名称、适用阶段和物理意义，而不是只写代数。
  \item 能在综合例题中标出力学、电磁、光学、热学和近代物理各自承担的模型角色。
  \item 能为后续学习写出一周复盘安排，并按错因类型整理错题。
\end{enumerate}
\section*{二、综合题的分段与统筹}

学完 30 节后，最重要的变化不是“记住了多少公式”，而是看到题目时能主动建立模型。稳定的解题能力来自先看清对象、过程和约束，再选择合适规律。

\textbf{科学发展视角：} 继续学习物理时，目标不是把公式越背越多，而是逐步形成“建模、计算、验证、修正”的科学工作方式。无论是高考综合题、竞赛入门题，还是未来接触人工智能、材料、能源、航天和医学工程，核心能力都是把真实问题转化为可分析的物理模型。

\begin{center}
% TikZ 图：综合题拆解流程
% Day 30：综合训练
\begin{tikzpicture}[scale=0.86,>=Latex]
  \node[draw,rounded corners=2pt,fill=blue!8,minimum width=1.55cm,minimum height=0.72cm] (a) at (0,0) {画图};
  \node[draw,rounded corners=2pt,fill=blue!8,minimum width=1.55cm,minimum height=0.72cm] (b) at (2.2,0) {选对象};
  \node[draw,rounded corners=2pt,fill=blue!8,minimum width=1.55cm,minimum height=0.72cm] (c) at (4.4,0) {分过程};
  \node[draw,rounded corners=2pt,fill=blue!8,minimum width=1.55cm,minimum height=0.72cm] (d) at (6.6,0) {列方程};
  \node[draw,rounded corners=2pt,fill=blue!8,minimum width=1.75cm,minimum height=0.72cm] (e) at (8.8,0) {复核结果};
  \foreach \p/\q in {a/b,b/c,c/d,d/e}{
    \draw[->,thick] (\p) -- (\q);
  }
  \node[font=\small,align=center] at (4.4,-0.9)
    {每一步都要写出物理含义，不能只堆公式。};
\end{tikzpicture}


\figcaption{综合题五步拆解：画图、选对象、分过程、列方程、复核结果}\label{fig:comprehensive-solving-flow}
\end{center}

\section*{三、综合解题主线}

\begin{conceptbox}{综合题五步法}

\begin{enumerate}
  \item \textbf{画图：} 把物体、轨道、场区、电路、方向都画出来。
  \item \textbf{选对象：} 明确研究单个物体还是系统。
  \item \textbf{分过程：} 找到力、电场、磁场、接触、碰撞等条件改变的位置。
  \item \textbf{列方程：} 每个方程旁边写出所用物理规律。
  \item \textbf{复核结果：} 检查单位、正负、极限和数量级。
\end{enumerate}

\textbf{提醒：} 综合题的难点通常不在单个公式，而在“什么时候换模型”。

\textbf{答题纸三栏法：} 草稿可以灵活，正式书写要让阅卷者看见建模过程。每一段建议按三栏写：
\begin{itemize}
  \item \textbf{状态或阶段：} 例如“斜面下滑”“碰撞瞬间”“进入磁场后”。
  \item \textbf{使用规律：} 例如“机械能守恒”“水平方向动量守恒”“洛伦兹力提供向心力”。
  \item \textbf{连接量：} 例如“碰前速度”“碰后共同速度”“进入磁场速度”。
\end{itemize}

\textbf{认知提醒：} 综合题不是把所有公式写在一起，而是让每个方程都有自己的阶段和物理理由。若说不清某个方程属于哪一段，通常说明过程还没有分清。
\end{conceptbox}

\begin{conceptbox}{结果检验的三把尺}

\begin{itemize}
  \item \textbf{量纲：} 表达式单位是否等于所求量单位？
  \item \textbf{极限：} 参数变大、变小或趋零时，结果是否合理？
  \item \textbf{方向：} 矢量结果是否符合受力、场力和运动趋势？
\end{itemize}

例如得到速度表达式 $v=\sqrt{2gh}$，量纲正确；当 $h$ 增大时速度增大，极限合理；方向还要由具体运动决定。
\end{conceptbox}

\begin{samepage}
\section*{四、守恒关系与适用边界}

\textbf{综合题审题前检查：}
\nopagebreak[4]
\begin{enumerate}
  \item 先按时间顺序分段：初态、过程、转折点、末态分别写清。
  \item 每一段先确定一条主线，常见主线是牛顿定律与运动学、守恒关系、图像或电路；必要时再补充约束关系，避免把互不相干的公式同时堆上去。
  \item 遇到跨章节情境时，先把物理模型拆开，再把共享量写在两段之间，例如出口速度、位移、电流、磁通量变化量。
  \item 每写一个方程，就在旁边标明适用阶段和条件，防止把上一段的速度、加速度或电压误用到下一段。
  \item 最后检查单位、方向、正负号和极端情况，确认结果有物理意义。
\end{enumerate}
\end{samepage}

\begin{center}
% TikZ 图：力电综合示意
% Day 30：综合训练
\begin{tikzpicture}[scale=0.88,>=Latex]
  % 轨道在场区边界处结束；进入场区后粒子不再受轨道约束。
  \fill[blue!8] (5.2,0.3) rectangle (8.2,2.9);
  \draw[thick] (0,0) -- (3.2,0)
    .. controls (4.2,0) and (4.2,1.6) .. (5.2,1.6);
  \foreach \x in {5.6,6.2,6.8,7.4}{
    \draw[->,blue,thick] (\x,2.65) -- (\x,1.95);
  }
  \fill[red] (0.7,0.15) circle (2.5pt) node[above] {$+q$};
  \draw[->,red,thick] (0.7,0.15) -- (1.55,0.15) node[midway,above] {$v_0$};
  \fill[red] (5.2,1.6) circle (2.2pt);
  \draw[->,red,thick] (5.2,1.6) -- (6.05,1.6)
    node[midway,above] {$v$};
  \draw[red,thick,dashed,domain=5.2:8.0,samples=60]
    plot (\x,{1.6-0.10*(\x-5.2)*(\x-5.2)});
  \node at (6.7,3.2) {匀强电场区};
  \node[font=\small,align=center] at (4.1,-0.65)
    {轨道段先用能量或牛顿定律；离轨进入电场后分解运动。};
\end{tikzpicture}


\figcaption{力电综合示意：轨道运动、电场偏转和能量方法常组合出现}\label{fig:integrated-track-field}
\end{center}

这类题通常包含多个阶段：
\begin{enumerate}
  \item 在轨道上运动：先判断是否光滑、是否有外力做功；光滑段常用机械能守恒，粗糙或有外力做功时用动能定理或功能关系。
  \item 离开轨道进入电场：先判断电场力方向，再把初速度沿电场方向和垂直电场方向分解。
  \item 离开场区：用速度分量合成或继续做匀速直线运动。
\end{enumerate}

\section*{五、跨模块综合模型}

\subsection*{模型 1：力学综合}

常见组合：
\[
\text{受力分析}\rightarrow\text{加速度}\rightarrow\text{运动学}
\]
或
\[
\text{高度/弹簧/摩擦}\rightarrow\text{能量方程}.
\]

若题中出现“最大压缩量、最远距离、刚好通过最高点”，往往要优先考虑能量和临界条件。

\subsection*{模型 2：碰撞与后续运动}

碰撞时间很短且所选系统在碰撞方向上的外力冲量可忽略时，碰撞阶段可用动量守恒；碰撞后的运动再根据轨道、摩擦或弹簧选择能量方法。不要把整个过程一口气套机械能守恒。

\subsection*{模型 3：电场中的粒子运动}

先加速后偏转是典型模型。若用 $|q|$ 表示电荷量大小，且 $U_1$ 是使粒子加速的电压大小，则
\[
|q|U_1=\frac12mv_0^2,
\]
进入偏转电场后，
\[
a=\frac{|q|E}{m}=\frac{|q|U_2}{md}.
\]
若只求偏移大小可用上式；实际偏转方向还要由电荷正负和电场方向判断。水平方向匀速，竖直方向匀加速。

\subsection*{模型 4：电磁感应中的动力学}

在闭合回路中，当导体棒、速度和匀强磁场方向两两垂直时，速度决定感应电动势，电流决定安培力，安培力又反过来影响速度：
\[
v\rightarrow \mathcal E=BLv\rightarrow I=\frac{\mathcal E}{R}\rightarrow F_A=BIL.
\]
这就是反馈模型，常出现“最终匀速”或“终端速度”。

\begin{center}
% TikZ 图：电磁感应综合题
% Day 30：综合训练
\begin{tikzpicture}[scale=0.9,>=Latex]
  \draw[thick] (-2.5,1.0) -- (2.8,1.0);
  \draw[thick] (-2.5,-1.0) -- (2.8,-1.0);
  \draw[thick] (-2.5,1.0) -- (-2.5,0.32);
  \draw[thick,fill=white] (-2.72,0.32) rectangle (-2.28,-0.32);
  \node[font=\small] at (-2.50,0) {$R$};
  \draw[thick] (-2.5,-0.32) -- (-2.5,-1.0);
  \draw[very thick,red] (-0.3,-1.0) -- (-0.3,1.0);
  \foreach \x in {-1.8,-1.0,-0.2,0.6,1.4,2.2} {
    \foreach \y in {-0.55,0.25} {
      \draw[blue] (\x-0.05,\y-0.05) -- (\x+0.05,\y+0.05);
      \draw[blue] (\x-0.05,\y+0.05) -- (\x+0.05,\y-0.05);
    }
  }
  \draw[->,purple,thick] (-0.3,0.45) -- (1.05,0.45) node[right] {$F$};
  \draw[->,red,thick] (-0.3,0.05) -- (0.85,0.05) node[right] {$v$};
  \draw[->,orange,thick] (-0.3,-0.35) -- (-1.35,-0.35) node[left] {$F_A$};
  \draw[->,blue!70!black,thick] (-0.3,0.62) -- (-0.3,0.93)
    node[right,font=\small] {$I$};
  \node[blue] at (3.15,0.8) {$B$ 向里};
  \node[font=\small,align=center] at (0,-1.65)
    {外力做功 $\rightarrow$ 电能 $\rightarrow$ 焦耳热；终端速度由 $F=F_A$ 决定};
\end{tikzpicture}


\figcaption{电磁感应综合：速度、电流、安培力和能量转化相互联系}\label{fig:comprehensive-em-induction}
\end{center}

\subsection*{模型 5：实验综合}

实验综合题通常不只问“算一个数”，而是按完整科学流程考查：
\[
\text{原理公式}\rightarrow\text{测量量}\rightarrow\text{图像处理}\rightarrow\text{误差分析}.
\]
若题目给数据表，优先想能否线性化；若题目给电路，先判断电表实际测到的电压和电流；若题目问改进方案，先指出误差来源以及它怎样影响结果，再说明改动如何减小系统误差或偶然误差。

\begin{center}
% TikZ 图：实验模型复盘
% Day 29：把实验题转化为测量链条
\begin{tikzpicture}[scale=0.92,>=Latex]
  \tikzset{
    card/.style={draw,rounded corners=4pt,thick,fill=gray!6,
      minimum width=3.05cm,minimum height=2.48cm,
      text width=2.68cm,align=center,font=\small},
    arrow/.style={->,thick},
    note/.style={font=\scriptsize,align=center}
  }

  \node[card] (direct) at (0,0) {
    \textbf{直接测量}\\[3pt]
    长度、时间\\
    电压、电流\\
    体积、面积
  };
  \node[card] (model) at (3.85,0) {
    \textbf{模型关系}\\[3pt]
    $v\approx \ell/\Delta t$\\
    $a\approx\Delta v/\Delta t$\\
    $U=E-Ir$\\
    $d\approx V/S$
  };
  \node[card] (target) at (7.70,0) {
    \textbf{间接量}\\[3pt]
    速度、加速度\\
    电动势、内阻\\
    分子大小
  };

  \draw[arrow] (direct) -- (model);
  \draw[arrow] (model) -- (target);

  \node[note,orange!85!black] at (3.85,-1.88)
    {读数后还要检查单位、有效数字和误差方向};
\end{tikzpicture}


\figcaption{实验综合的共同结构：直接读数要通过模型关系转化为待求量}\label{fig:experiment-model-review-day30}
\end{center}

\textbf{典型转化：}
\begin{itemize}
  \item 纸带和光电门：把长度、时间转化为速度或加速度。
  \item 伏安法、电源实验、电桥：把电压、电流或平衡关系转化为电阻、电动势和内阻。
  \item 油膜法：把微小体积和可测面积转化为分子大小的数量级。
  \item 图像法：斜率、截距和面积都必须带单位，不能只看几何形状。
\end{itemize}

\subsection*{模型 6：后段专题综合}

光学、热学和近代物理通常不再像力学那样直接从“受力列方程”开始，但解题方式仍然是先建模。光学先画光路，判断反射、折射、全反射、色散或波动光学；热学先区分状态量和过程量，再用第一定律和准静态过程在 $p$-$V$ 图中图线下的有向面积；近代物理先判断粒子、光子、能级或核反应，再用能量、动量、电荷数和质量数守恒。

\begin{center}
% TikZ 图：Day 30 后段专题综合判断
\begin{tikzpicture}[scale=0.92,>=Latex]
  \tikzset{
    box/.style={draw,rounded corners=3pt,thick,fill=gray!6,
      minimum width=3.05cm,minimum height=2.35cm,
      text width=2.62cm,align=center,font=\small},
    title/.style={font=\bfseries},
    ray/.style={->,thick,blue},
    level/.style={thick},
    note/.style={font=\scriptsize,align=center}
  }

  \node[box] (opt) at (0,0) {};
  \node[title] at (0,0.82) {光学};
  \draw[thick] (-1.05,0.18) -- (1.05,0.18);
  \draw[ray] (-1.05,-0.45) -- (-0.20,0.18);
  \draw[ray] (-0.20,0.18) -- (0.82,0.66);
  \draw[dashed,blue] (-0.20,0.18) -- (0.82,-0.30);
  \node[note] at (0,-0.87) {先画光路\\再判反射、折射或成像};

  \node[box] (therm) at (4.0,0) {};
  \node[title] at (4.0,0.82) {热学};
  \draw[->,thick] (3.0,-0.55) -- (3.0,0.55) node[left] {$p$};
  \draw[->,thick] (3.0,-0.55) -- (4.95,-0.55) node[below] {$V$};
  \draw[very thick,orange] (3.18,-0.35) -- (4.55,0.38);
  \fill[orange!20] (3.18,-0.55) -- (3.18,-0.35) -- (4.55,0.38) -- (4.55,-0.55) -- cycle;
  \node[note] at (4.0,-0.87) {区分状态量和过程量\\面积对应气体做功};

  \node[box] (modern) at (8.0,0) {};
  \node[title] at (8.0,0.82) {近代物理};
  \draw[level] (7.15,0.35) -- (8.85,0.35) node[right] {$E_3$};
  \draw[level] (7.15,-0.05) -- (8.85,-0.05) node[right] {$E_2$};
  \draw[level] (7.15,-0.48) -- (8.85,-0.48) node[right] {$E_1$};
  \draw[->,red,thick] (7.58,0.30) -- (7.58,-0.42);
  \node[red,font=\scriptsize] at (7.15,-0.05) {$h\nu$};
  \node[note] at (8.0,-0.87) {先看能量量子化\\再用守恒关系};
\end{tikzpicture}


\figcaption{后段专题综合：光路、热过程和能级图也要先转化为模型}\label{fig:comprehensive-late-topics}
\end{center}

\textbf{判断顺序：}
\begin{enumerate}
  \item 光学题先问“光经过哪些界面、是否发生全反射或色散、是否需要判断相干和光程差”。
  \item 热学题先问“状态是否改变、系统吸热还是放热、气体对外做功还是外界对气体做功”，并统一符号约定。
  \item 近代物理题先问“能量是否量子化、守恒量有哪些”。
\end{enumerate}

\subsection*{模型 7：能源与可持续发展}

高中物理还要训练一种更大的建模视角：真实技术系统不是一个孤立公式，而是一条能量链。分析能源与工程问题时，先问能量从哪里来、经过哪些转换、在哪里传输和储存、最后变成什么有用效果，又在哪里损耗。

\begin{center}
% TikZ 图：能源与可持续发展系统模型
% Day 30：综合收束
\begin{tikzpicture}[>=Latex,scale=0.88,every node/.style={font=\scriptsize}]
  \tikzset{
    block/.style={draw,rounded corners=2pt,minimum width=2.45cm,minimum height=1.05cm,align=center,fill=gray!6},
    note/.style={draw,rounded corners=2pt,align=center,fill=yellow!12,inner sep=4pt},
    flow/.style={->,thick},
    loss/.style={->,thick,red!75},
    info/.style={->,thick,blue!75}
  }

  \node[block,fill=green!10] (source) at (0,1.6) {一次能源\\太阳能、风能\\化石燃料、核能};
  \node[block,fill=orange!10] (convert) at (3.2,1.6) {能量转换\\发电机、热机\\光伏电池};
  \node[block,fill=blue!8] (grid) at (6.4,1.6) {传输与储能\\升压输电\\电池、超级电容};
  \node[block,fill=purple!8] (use) at (9.6,1.6) {终端利用\\电动机、照明\\计算与通信};

  \draw[flow] (source) -- (convert);
  \draw[flow] (convert) -- (grid);
  \draw[flow] (grid) -- (use);

  \node[note] (law) at (3.2,-0.15) {守恒账本\\输入能量 $=$ 有用能\\$+$ 损耗能};
  \node[note] (eta) at (6.4,-0.15) {效率意识\\$\eta=\dfrac{E_{\text{有用}}}{E_{\text{输入}}}$\\高压输电减小 $I^2R$};
  \node[note] (risk) at (9.6,-0.15) {边界意识\\安全、环境\\成本与可靠性};

  \draw[info] (convert.south) -- (law.north);
  \draw[info] (grid.south) -- (eta.north);
  \draw[info] (use.south) -- (risk.north);

  \node[align=center,font=\tiny] at (4.8,2.85) {物理不是只算一个环节，而是检查整条能量链};
\end{tikzpicture}


\figcaption{能源与可持续发展模型：从能量来源、转换、传输到利用都要检查守恒、效率和边界}\label{fig:energy-sustainable-system}
\end{center}

\textbf{三个判断点：}
\begin{enumerate}
  \item \textbf{守恒条件清单：} 总能量守恒，但能量经过实际转换后，可用于完成目标任务的比例通常降低；废热、噪声和线路损耗都会减少有用能量。
  \item \textbf{效率意识：} 热机、电动机、输电线、变压器和充电器都有效率问题；功率大不等于效率高。
  \item \textbf{边界意识：} 核能、化石能源、可再生能源和储能技术都要同时考虑安全、环境、成本、稳定性和材料条件。
\end{enumerate}

\subsection*{模型 8：数量级估算}

真实问题常常不给出所有精确数据，或者数据太多、太杂。此时先做数量级估算，可以判断答案是否合理，也能帮助发现漏单位、漏平方、指数写错等问题。数量级估算不是任意猜测，而是有步骤的近似建模。

\begin{center}
% TikZ 图：数量级估算模型
% Day 30：综合收束
\begin{tikzpicture}[>=Latex,scale=0.9,every node/.style={font=\scriptsize}]
  \tikzset{
    estep/.style={draw,rounded corners=2pt,minimum width=2.05cm,minimum height=0.86cm,align=center,fill=blue!7},
    check/.style={draw,rounded corners=2pt,minimum width=2.35cm,minimum height=0.86cm,align=center,fill=orange!10},
    arrow/.style={->,thick}
  }

  \node[estep] (scene) at (0,0) {真实情境\\先画草图};
  \node[estep] (main) at (2.65,0) {抓主要量\\忽略次要细节};
  \node[estep] (value) at (5.3,0) {取近似值\\一到两位有效数字};
  \node[estep] (calc) at (7.95,0) {代入估算\\保留量级};
  \node[check] (test) at (10.75,0) {检查结果\\单位、极限、常识};

  \foreach \a/\b in {scene/main,main/value,value/calc,calc/test}{
    \draw[arrow] (\a) -- (\b);
  }

  \node[draw,rounded corners=2pt,fill=gray!6,align=left,inner sep=5pt] at (5.35,-1.55)
    {例：估算一天用电量，不必先知道每个电器的精确功率，\\
     先列主要用电器，再用 $E=Pt$ 判断数量级是否合理。};
\end{tikzpicture}


\figcaption{数量级估算：先抓主要量，再用单位、极限和常识检查结果}\label{fig:order-of-magnitude-estimate}
\end{center}

\textbf{估算原则：}
\begin{enumerate}
  \item \textbf{先定主量：} 找出最影响结果的物理量，如功率、时间、速度、面积或距离。
  \item \textbf{再取近似：} 只保留一到两位有效数字，优先判断 $10^0,10^1,10^2$ 这样的量级。
  \item \textbf{最后检验：} 单位必须正确，结果不能明显违背生活经验或物理边界。
\end{enumerate}

\textbf{例：} 估算一盏 $10\ \mathrm{W}$ 的 LED 灯连续亮 $5\ \mathrm{h}$ 消耗的电能：
\[
E=Pt\approx 10\ \mathrm{W}\times 5\ \mathrm{h}=50\ \mathrm{Wh}=0.05\ \mathrm{kWh}.
\]
如果算出 $50\ \mathrm{kWh}$，就应检查数量级：这已经接近很多家庭数天甚至更久的用电量，明显不合理。

\newpage
\section*{六、例题精讲}

\begin{examplebox}{力学综合：弹簧与粗糙面}
\textbf{题意分析：} 质量为 $m$ 的物块在水平粗糙面上运动，刚接触劲度系数为 $k$ 的轻弹簧时速度为 $v_0$。物块继续向右滑动，将弹簧压缩至最大量 $x$；物块与水平面间的动摩擦因数为 $\mu$。求 $x$ 满足的方程。

\textbf{物理模型：} 动能转化为弹性势能，同时摩擦力做负功。

\begin{center}
% TikZ 图：粗糙水平面上压缩弹簧
% Day 30：综合训练
\begin{tikzpicture}[scale=0.88,>=Latex,line cap=round,line join=round]
  \path[use as bounding box] (-0.35,-0.82) rectangle (7.80,2.08);
  \draw[thick] (0,0) -- (7.35,0);
  \foreach \x in {0.25,0.75,...,7.0}{
    \draw[gray] (\x,0) -- ++(0.18,-0.18);
  }
  \draw[thick] (7.35,0) -- (7.35,1.35);
  \draw[decorate,decoration={coil,aspect=0.45,segment length=4pt,amplitude=4pt},thick]
    (5.35,0.55) -- (7.35,0.55);

  % 刚接触弹簧时的位置（虚线）与最大压缩位置（实线）分开画。
  \draw[dashed,gray!80,fill=gray!6] (3.35,0.18) rectangle (4.15,0.92);
  \node[gray!80] at (3.75,0.55) {$m$};
  \draw[->,red,thick] (3.35,1.16) -- (4.15,1.16)
    node[midway,above] {$v_0$};

  \draw[fill=gray!14,thick] (4.55,0.18) rectangle (5.35,0.92);
  \node at (4.95,0.55) {$m$};
  \draw[->,orange,thick] (5.25,1.43) -- (4.35,1.43)
    node[midway,above] {$f=\mu mg$};

  \draw[dashed,blue!65] (4.15,0.08) -- (4.15,-0.35);
  \draw[dashed,blue!65] (5.35,0.08) -- (5.35,-0.35);
  \draw[<->,blue,thick] (4.15,-0.42) -- (5.35,-0.42)
    node[midway,below] {$x$};
  \node[font=\scriptsize,align=center] at (2.08,1.55)
    {虚线：刚接触弹簧};
  \node[font=\scriptsize,align=center] at (6.62,1.58)
    {最大压缩\\$v=0$};
\end{tikzpicture}

\end{center}

\textbf{解题过程：}
从接触弹簧到最大压缩，初态动能为 $\dfrac12mv_0^2$，末态速度为 0，弹簧势能为 $\dfrac12kx^2$，摩擦力做功为 $-\mu mgx$。
\[
\frac12mv_0^2-\mu mgx=\frac12kx^2.
\]
整理：
\[
\frac12kx^2+\mu mgx-\frac12mv_0^2=0.
\]
取满足 $x>0$ 的物理解：
\[
\boxed{x=\frac{\sqrt{\mu^2m^2g^2+kmv_0^2}-\mu mg}{k}}.
\]

\textbf{方法总结：} 综合能量题先检查是否有摩擦耗散。有滑动摩擦时机械能通常不守恒，但动能定理或功能关系仍可用。
\end{examplebox}

\begin{examplebox}{电场综合：先加速后偏转}
\textbf{题意分析：} 质量为 $m$、电荷量为 $q$ 的粒子由静止经电压大小为 $U_1$ 的电场加速后，以水平速度从中线进入板长为 $L$、板间距为 $d$、电压大小为 $U_2$ 的匀强偏转电场。忽略重力，并假定粒子能从极板右端飞出，求离开偏转电场时竖直偏移量的大小。

\textbf{物理模型：} 加速区用动能定理，偏转区用类平抛运动。

\begin{center}
% TikZ 图：先加速后偏转模型
% Day 17：带电粒子在电场中的运动
\begin{tikzpicture}[scale=0.78,>=Latex]
  % 加速区
  \draw[thick] (-3.2,1.0) -- (-3.2,-1.0);
  \draw[thick] (-1.6,1.0) -- (-1.6,-1.0);
  \node[above] at (-3.2,1.04) {$+$};
  \node[above] at (-1.6,1.04) {$-$};
  \foreach \y in {-0.45,0,0.45} {
    \draw[->,blue] (-3.0,\y) -- (-1.8,\y);
  }
  \draw[<->] (-3.2,1.32) -- (-1.6,1.32);
  \node[above,blue,fill=white,inner sep=1pt] at (-2.4,1.32) {$U_1$};
  \node[below left,red,fill=white,inner sep=1pt] at (-3.6,0) {$+q$};
  \draw[red,thick,dashed] (-3.6,0) -- (-1.62,0);
  \draw[-{Stealth[length=5pt,width=6pt]},red,line width=1.05pt] (-1.35,0) -- (-0.35,0);
  \node[red,fill=white,inner sep=1pt] at (-0.82,0.28) {$v_0$};

  % 偏转区
  \draw[thick] (-0.2,1.2) -- (3.0,1.2);
  \draw[thick] (-0.2,-1.2) -- (3.0,-1.2);
  \node[above] at (-0.35,1.2) {$+$};
  \node[below] at (-0.35,-1.2) {$-$};
  \foreach \x in {0.35,1.05,1.75,2.45} {
    \draw[->,blue] (\x,1.0) -- (\x,-1.0);
  }
  \draw[<->] (-0.2,1.55) -- (3.0,1.55);
  \node[above,fill=white,inner sep=1pt] at (1.4,1.55) {$L$};
  \draw[<->] (3.35,1.2) -- (3.35,-1.2);
  \node[right,fill=white,inner sep=1pt] at (3.35,0) {$d$};
  \node[blue,fill=white,inner sep=1pt] at (0.15,-1.48) {$U_2$};

  \draw[red,thick,dashed,domain=0.25:3.3,samples=80]
    plot (\x,{-0.10*(\x+0.2)*(\x+0.2)});
  \draw[dashed,gray] (3.3,0) -- (3.3,-1.23);
  \draw[<->] (4.0,0) -- (4.0,-1.23);
  \node[right,fill=white,inner sep=1pt] at (4.0,-0.62) {$y$};

  \node[font=\small,fill=white,inner sep=1pt] at (-2.4,-1.48) {加速区};
  \node[font=\small,fill=white,inner sep=1pt] at (2.15,-1.82) {偏转区};
\end{tikzpicture}

\end{center}

\textbf{解题过程：}
加速区：
\[
|q|U_1=\frac12mv_0^2,
\qquad
v_0=\sqrt{\frac{2|q|U_1}{m}}.
\]
偏转区电场强度
\[
E=\frac{U_2}{d},
\qquad
a=\frac{|q|E}{m}=\frac{|q|U_2}{md}.
\]
通过偏转区时间
\[
t=\frac{L}{v_0}.
\]
竖直偏移
\[
y=\frac12at^2
=\frac12\cdot\frac{|q|U_2}{md}\cdot\frac{L^2}{v_0^2}.
\]
代入 $v_0^2=\dfrac{2|q|U_1}{m}$：
\[
y=\frac{U_2L^2}{4U_1d}.
\]

\textbf{飞出条件：} 因粒子从中线进入，要从两板之间飞出，还须满足
\[
y<\frac d2,
\]
即上述结果必须小于半板间距；否则粒子会先打到极板，不能再使用“运动时间为 $L/v_0$”来描述离场状态。

\textbf{方法总结：} 上式求的是偏移大小。结果中 $m$、$q$ 消去，说明在忽略重力、初速由同一加速电压产生且粒子能飞出极板时，偏移大小由电压和几何尺寸决定；偏转向上还是向下仍要看电荷正负和电场方向。
\end{examplebox}

\begin{examplebox}{磁场综合：电场加速后进入磁场}
\textbf{题意分析：} 带电粒子质量为 $m$、电荷量为 $q$，从静止经电压大小为 $U$ 的电场加速后，垂直进入磁感应强度为 $B$ 的匀强磁场。忽略相对论效应，求粒子在磁场中的轨道半径和周期。

\textbf{物理模型：} 加速区用动能定理；磁场区用洛伦兹力提供向心力。

\begin{center}
% TikZ 图：电场加速后进入磁场
% Day 29：模型复盘
\begin{tikzpicture}[scale=0.84,>=Latex,line cap=round,line join=round]
  \path[use as bounding box] (-0.35,-0.95) rectangle (8.10,1.75);
  \draw[thick] (0.30,-0.45) -- (0.30,1.05);
  \draw[thick] (2.05,-0.45) -- (2.05,1.05);
  \node[above] at (0.30,1.05) {$+$};
  \node[above] at (2.05,1.05) {$-$};
  \node[above] at (1.18,1.05) {$U$};
  \foreach \y in {-0.18,0.25,0.68}{
    \draw[->,blue,thick] (0.58,\y) -- (1.72,\y);
  }
  % Place the initially slow positive charge inside the accelerating gap.
  % Putting it outside the positive plate would leave it outside the
  % idealized uniform field shown here.
  \fill[red] (0.58,0.25) circle (1.8pt);
  \node[red,below] at (0.58,0.20) {$+q$};
  \draw[->,red,thick,dashed] (0.72,0.25) -- (2.65,0.25);
  \node[red,font=\scriptsize,above] at (2.25,0.25) {$v_0$};
  \draw[thick] (2.55,-0.55) rectangle (7.55,1.30);
  \foreach \x in {3.0,3.7,4.4,5.1,5.8,6.5,7.2}{
    \foreach \y in {-0.25,0.58,1.15}{
      \node[blue] at (\x,\y) {$\times$};
    }
  }
  % For +q moving right in a field into the page, the center lies directly
  % above the entry point.  The red trajectory and radius share this center.
  \coordinate (P) at (2.65,0.25);
  \coordinate (C) at (2.65,0.95);
  \draw[red,very thick,->] (P) arc[start angle=-90,end angle=0,radius=0.70];
  \fill (C) circle (1.3pt);
  \node[font=\scriptsize,anchor=east] at (2.54,0.98) {$C$};
  \draw[orange!85!black,thick] (P) -- (C);
  \node[orange!85!black,font=\scriptsize,anchor=east] at (2.54,0.60) {$r$};
  \node[font=\small] at (1.18,1.48) {电场加速};
  \node[font=\small] at (5.25,1.48) {磁场偏转};
  \node[blue,font=\scriptsize] at (6.40,-0.30) {$\vec B$ 垂直纸面向里};
\end{tikzpicture}

\end{center}

\textbf{解题过程：}
加速区：
\[
|q|U=\frac12mv^2,\qquad v=\sqrt{\frac{2|q|U}{m}}.
\]
磁场区：
\[
|q|vB=\frac{mv^2}{r},\qquad r=\frac{mv}{|q|B}.
\]
代入速度：
\[
r=\frac{1}{B}\sqrt{\frac{2mU}{|q|}}.
\]
周期为
\[
T=\frac{2\pi m}{|q|B}.
\]

\textbf{方法总结：} 这里用 $|q|$ 表示电荷量大小，且加速电场对粒子做正功。图示以正电荷为例；若粒子带负电，进入磁场后的偏转方向相反，但半径和周期大小不变。本题中电场力做功改变动能，磁场力提供向心力并改变速度方向。看到“先加速后进入磁场”，应先把加速段和磁场偏转段分开处理。
\end{examplebox}

\begin{examplebox}{电磁感应综合：恒力拉导体棒}
\textbf{题意分析：} 长为 $L$ 的导体棒在间距为 $L$ 的光滑导轨上向右运动，并始终与导轨良好接触。匀强磁场 $B$ 垂直导轨平面，闭合回路的总电阻恒为 $R$。水平向右的恒力 $F$ 作用于棒，忽略自感，求导体棒的终端速度，并说明终端状态下的能量转化。

\textbf{物理模型：} 导体棒运动产生感应电动势，感应电流产生反向安培力，最终受力平衡。

\begin{center}
% TikZ 图：电磁感应综合题
% Day 30：综合训练
\begin{tikzpicture}[scale=0.9,>=Latex]
  \draw[thick] (-2.5,1.0) -- (2.8,1.0);
  \draw[thick] (-2.5,-1.0) -- (2.8,-1.0);
  \draw[thick] (-2.5,1.0) -- (-2.5,0.32);
  \draw[thick,fill=white] (-2.72,0.32) rectangle (-2.28,-0.32);
  \node[font=\small] at (-2.50,0) {$R$};
  \draw[thick] (-2.5,-0.32) -- (-2.5,-1.0);
  \draw[very thick,red] (-0.3,-1.0) -- (-0.3,1.0);
  \foreach \x in {-1.8,-1.0,-0.2,0.6,1.4,2.2} {
    \foreach \y in {-0.55,0.25} {
      \draw[blue] (\x-0.05,\y-0.05) -- (\x+0.05,\y+0.05);
      \draw[blue] (\x-0.05,\y+0.05) -- (\x+0.05,\y-0.05);
    }
  }
  \draw[->,purple,thick] (-0.3,0.45) -- (1.05,0.45) node[right] {$F$};
  \draw[->,red,thick] (-0.3,0.05) -- (0.85,0.05) node[right] {$v$};
  \draw[->,orange,thick] (-0.3,-0.35) -- (-1.35,-0.35) node[left] {$F_A$};
  \draw[->,blue!70!black,thick] (-0.3,0.62) -- (-0.3,0.93)
    node[right,font=\small] {$I$};
  \node[blue] at (3.15,0.8) {$B$ 向里};
  \node[font=\small,align=center] at (0,-1.65)
    {外力做功 $\rightarrow$ 电能 $\rightarrow$ 焦耳热；终端速度由 $F=F_A$ 决定};
\end{tikzpicture}

\end{center}

\textbf{解题过程：}
\[
\mathcal E=BLv,\qquad I=\frac{\mathcal E}{R}=\frac{BLv}{R}.
\]
安培力大小
\[
F_A=BIL=\frac{B^2L^2v}{R}.
\]
最终匀速时
\[
F=F_A,
\]
所以
\[
v=\frac{FR}{B^2L^2}.
\]
能量上，外力做功的功率为
\[
P_{\text{外}}=Fv,
\]
在回路中只有电阻发热、忽略机械摩擦和辐射等损耗的理想模型中，最终匀速时动能不再变化，外力输入功率转化为电路焦耳热：
\[
P_R=I^2R.
\]

\textbf{方法总结：} 电磁感应综合题先画反馈链：运动产生电，电流产生力，力再限制运动。这里的“终端速度”是受力逐渐趋于平衡时的极限速度，并不表示棒在此前始终做匀加速运动。
\end{examplebox}

\begin{examplebox}{实验综合：由图像求电源内阻}
\textbf{题意分析：} 测电源电动势和内阻实验中，作出 $U$-$I$ 图像，纵截距为 $1.50\ \mathrm{V}$，图线斜率为 $-0.80\ \mathrm{V/A}$。求电动势和内阻，并说明为什么斜率取绝对值。

\textbf{物理模型：} 闭合电路欧姆定律
\[
U=E-Ir.
\]

\begin{center}
% TikZ 图：U-I 图像求电源电动势和内阻
% Day 30：综合训练
\begin{tikzpicture}[scale=0.90,>=Latex,line cap=round,line join=round]
  \draw[->,thick] (0,0) -- (5.10,0) node[right] {$I$};
  \draw[->,thick] (0,0) -- (0,2.75) node[above] {$U$};
  \draw[blue,thick] (0,2.35) -- (4.35,1.25);
  \node[left] at (0,2.35) {$E$};
  \draw[dashed] (3.20,1.54) -- (3.20,0);
  \draw[dashed] (3.20,1.54) -- (0,1.54);
  \fill[blue] (0,2.35) circle (1.7pt);
  \fill[blue] (3.20,1.54) circle (1.7pt);
  \node[blue,font=\small,align=center] at (3.65,2.45)
    {斜率 $k=-r$\\纵截距 $E$};
  \draw[<->,red,thick] (0.18,2.28) -- (3.02,1.58)
    node[midway,below,sloped] {$\Delta U/\Delta I$};
\end{tikzpicture}

\end{center}

\textbf{解题过程：}
与一次函数 $y=b+kx$ 对比：
\[
U=E-rI.
\]
纵截距就是
\[
E=1.50\ \mathrm{V}.
\]
斜率为 $-r$，所以
\[
r=0.80\ \Omega.
\]

\textbf{方法总结：} 图像题先读截距、斜率和单位；$U$-$I$ 图斜率为负，是因为内阻电压降随电流增大。
\end{examplebox}

\begin{learningbox}{综合题完成前四问}
综合题算完以后不要立刻停笔，至少快速问四个问题：
\begin{enumerate}
  \item \textbf{阶段是否完整：} 每个条件改变的位置是否都单独处理了？
  \item \textbf{方程是否有来源：} 每个方程旁边能否说出物理规律和适用条件？
  \item \textbf{结果是否有单位：} 速度、力、能量、电流、电阻等单位是否与题目目标量一致？
  \item \textbf{极限是否合理：} 摩擦趋零、磁场变强、电压增大、半径变大时，结果变化是否符合直觉？
\end{enumerate}
这四问不保证答案一定正确，但能及时拦住大量“模型用错、漏阶段、单位错、方向错”的失误。
\end{learningbox}

\section*{七、变式训练}

\begin{enumerate}
  \item 物块从光滑斜面滑下进入粗糙水平面，求停止距离，应如何分阶段？
  \begin{center}
  \begin{tikzpicture}[scale=0.8,>=Stealth,line cap=round,line join=round]
  \coordinate (A) at (0,2.0);
  \coordinate (B) at (2.7,0);
  \coordinate (C) at (5.4,0);
  % 连续接触面画在物块下方，避免断线被误读为悬空斜面。
  \draw[thick] (A) -- (B) -- (C);
  \begin{scope}[shift={(0.92,1.32)},rotate=-36.5]
    \draw[fill=gray!12,thick] (-0.32,0) rectangle (0.32,0.42);
  \end{scope}
  \draw[fill=gray!12] (3.55,0) rectangle (4.15,0.45);
  \draw[->,blue,thick] (1.25,1.03) -- (1.85,0.58);
  \draw[->,red,thick] (4.15,0.25) -- (3.45,0.25) node[midway,above] {$f$};
  \node[align=center,font=\small] at (1.25,-0.72) {阶段 1\\光滑斜面};
  \node[align=center,font=\small] at (4.45,-0.72) {阶段 2\\粗糙水平面};
  \draw[dashed] (2.7,0) -- (2.7,1.0);
  \node[above,font=\small] at (2.7,1.0) {模型切换};
\end{tikzpicture}

  \end{center}
  \item 完全非弹性碰撞后连接弹簧，碰撞和压缩两个阶段分别用什么规律？
  \begin{center}
  \begin{tikzpicture}[scale=0.88,>=Stealth,line cap=round,line join=round]
  \path[use as bounding box] (-0.25,-0.75) rectangle (8.15,2.25);
  \begin{scope}
    \draw[thick] (0,0) -- (2.8,0);
    \draw[fill=gray!12,thick] (0.25,0) rectangle (0.85,0.55) node[midway] {$m_1$};
    \draw[fill=gray!12,thick] (1.75,0) rectangle (2.35,0.55) node[midway] {$m_2$};
    \draw[->,red,thick] (0.9,0.38) -- (1.55,0.38) node[midway,above,font=\small] {$v_0$};
    \node[font=\small] at (1.4,-0.42) {碰撞前};
  \end{scope}
  \draw[->,thick] (3.0,0.35) -- (3.75,0.35);
  \begin{scope}[xshift=4.0cm]
    \draw[thick] (0,0) -- (3.8,0);
    \draw[fill=gray!12,thick] (0.25,0) rectangle (1.35,0.55) node[midway] {$m_1+m_2$};
    \draw[decorate,decoration={zigzag,segment length=5pt,amplitude=2pt},thick]
      (1.35,0.28) -- (2.28,0.28);
    \draw[thick] (2.28,0) rectangle (2.42,0.85);
    \draw[thick] (2.42,0) -- (2.42,1.2);
    \draw[thick] (2.42,0) -- (2.95,0);
    \node[blue,font=\scriptsize] at (1.82,0.92) {弹簧受压};
    \node[font=\small] at (1.75,-0.42) {碰后压缩弹簧};
  \end{scope}
  \node[font=\small,align=center] at (4.0,2.0)
    {碰撞：动量守恒\\压缩：无耗散用机械能，有耗散用功能关系};
\end{tikzpicture}

  \end{center}
  \item 带电粒子在匀强磁场中做圆周运动，若速度变为原来 2 倍，半径怎样变化？
  \item 导体棒在磁场中最终匀速，安培力与外力关系如何？
  \begin{center}
  % TikZ 图：电磁感应综合题
% Day 30：综合训练
\begin{tikzpicture}[scale=0.9,>=Latex]
  \draw[thick] (-2.5,1.0) -- (2.8,1.0);
  \draw[thick] (-2.5,-1.0) -- (2.8,-1.0);
  \draw[thick] (-2.5,1.0) -- (-2.5,0.32);
  \draw[thick,fill=white] (-2.72,0.32) rectangle (-2.28,-0.32);
  \node[font=\small] at (-2.50,0) {$R$};
  \draw[thick] (-2.5,-0.32) -- (-2.5,-1.0);
  \draw[very thick,red] (-0.3,-1.0) -- (-0.3,1.0);
  \foreach \x in {-1.8,-1.0,-0.2,0.6,1.4,2.2} {
    \foreach \y in {-0.55,0.25} {
      \draw[blue] (\x-0.05,\y-0.05) -- (\x+0.05,\y+0.05);
      \draw[blue] (\x-0.05,\y+0.05) -- (\x+0.05,\y-0.05);
    }
  }
  \draw[->,purple,thick] (-0.3,0.45) -- (1.05,0.45) node[right] {$F$};
  \draw[->,red,thick] (-0.3,0.05) -- (0.85,0.05) node[right] {$v$};
  \draw[->,orange,thick] (-0.3,-0.35) -- (-1.35,-0.35) node[left] {$F_A$};
  \draw[->,blue!70!black,thick] (-0.3,0.62) -- (-0.3,0.93)
    node[right,font=\small] {$I$};
  \node[blue] at (3.15,0.8) {$B$ 向里};
  \node[font=\small,align=center] at (0,-1.65)
    {外力做功 $\rightarrow$ 电能 $\rightarrow$ 焦耳热；终端速度由 $F=F_A$ 决定};
\end{tikzpicture}

  \end{center}
  \item 对任意一个综合题，写出“选对象、分过程、列方程”的三句话。
\end{enumerate}

\begin{learningbox}{变式训练参考结果}
\begin{enumerate}[itemsep=0.45em,topsep=0.3em,parsep=0pt]
  \item 按接触面分为两段：光滑斜面段可用机械能守恒求进入水平面时的速度；粗糙水平面段用动能定理或匀减速运动求停止距离。
  \item 碰撞时间很短且碰撞方向上的外力冲量可忽略时，先用动量守恒求共同速度；随后压缩弹簧，若无摩擦用机械能守恒，有摩擦则用功能关系。
  \item 由 $r=mv/(|q|B)$，在 $m$、$|q|$、$B$ 不变时，速度变为原来的 2 倍，半径也变为原来的 2 倍。
  \item 最终匀速时合力为零，安培力与外力等大反向；若忽略其他耗散，外力做功转化为回路中的焦耳热。
  \item 示例表述：下滑阶段选“物块与地球系统”，碰撞阶段选“两物体系统”，压缩阶段选“两物体与弹簧系统”（涉及重力势能时再含地球）；按“下滑、碰撞、压缩弹簧”分过程，分别列机械能守恒、动量守恒和功能关系。
\end{enumerate}
\end{learningbox}

\section*{八、综合能力拓展}
\begin{competitionbox}
\textbf{后续学习路线}

\begin{center}
% TikZ 图：后续学习路线
% Day 30：学习路线
\begin{tikzpicture}[scale=0.86,>=Latex]
  \node[draw,rounded corners=2pt,fill=green!10,minimum width=1.6cm] (basic) at (0,0) {课内基础};
  \node[draw,rounded corners=2pt,fill=blue!8,minimum width=1.6cm] (model) at (3.0,0) {模型熟练};
  \node[draw,rounded corners=2pt,fill=orange!12,minimum width=1.8cm] (exp) at (6.1,0) {实验与图像};
  \node[draw,rounded corners=2pt,fill=red!8,minimum width=1.8cm] (comp) at (9.3,0) {综合与竞赛};
  \foreach \a/\b in {basic/model,model/exp,exp/comp}{\draw[->,thick] (\a) -- (\b);}
  \node[font=\small,align=center] at (0,-0.8) {定义\\单位\\方向};
  \node[font=\small,align=center] at (3.0,-0.8) {受力\\能量\\动量};
  \node[font=\small,align=center] at (6.1,-0.8) {数据\\误差\\斜率};
  \node[font=\small,align=center] at (9.3,-0.8) {分段\\近似\\推广};
\end{tikzpicture}


\figcaption{一个月入门后的学习路线：基础、模型、实验、综合逐步推进}\label{fig:learning-route-map}
\end{center}

\textbf{课内同步路线：}
\begin{enumerate}
  \item 每学一节，先写概念定义和适用条件。
  \item 每整理一道错题，标注错因：概念、受力、公式、计算、审题还是表达。
  \item 每周复盘一次模型，而不是只复盘题号。
\end{enumerate}

\textbf{竞赛入门路线：}
\begin{enumerate}
  \item 力学：加强矢量、微元、能量、动量和圆周运动。
  \item 电磁学：加强场的叠加、电势、带电粒子运动和电磁感应综合。
  \item 数学工具：熟悉函数图像、三角关系、近似思想和入门微积分语言。
\end{enumerate}

\textbf{最重要的习惯：} 每一个公式都要能回答三件事：从哪里来、适用于什么条件、结果变大变小是否合理。
\end{competitionbox}

\subsection*{竞赛与综合提升：复盘系统}

一个月学习结束后，最容易出现的问题是“课堂上理解了，但隔一段时间后模型又混在一起”。解决办法不是只靠增加题量，而是建立复盘系统。

\begin{center}
% TikZ 图：复盘循环
% Day 30：综合训练
\begin{tikzpicture}[scale=0.9,>=Latex]
  \node[draw,rounded corners=2pt,fill=green!10,minimum width=1.35cm] (a) at (0,2.1) {做题};
  \node[draw,rounded corners=2pt,fill=blue!8,minimum width=1.65cm] (b) at (4.0,2.1) {定位错因};
  \node[draw,rounded corners=2pt,fill=orange!12,minimum width=1.65cm] (c) at (4.0,0) {归入模型};
  \node[draw,rounded corners=2pt,fill=red!8,minimum width=1.65cm] (d) at (0,0) {重做变式};
  \draw[->,thick] (a) -- (b);
  \draw[->,thick] (b) -- (c);
  \draw[->,thick] (c) -- (d);
  \draw[->,thick] (d) -- (a);
  \node[align=center,font=\small] at (2.0,1.05)
    {错题不是收藏品，\\要变成模型经验};
\end{tikzpicture}


\figcaption{复盘循环：把错题转化为模型经验}\label{fig:comprehensive-review-cycle}
\end{center}

\subsection*{1.\quad 错题分类表}

\begin{center}
\begin{tabular}{lll}
\toprule
错因类型 & 典型表现 & 处理方式 \\
\midrule
概念错误 & 把热量当状态量、把电势和电势能混淆 & 回到定义和适用条件 \\
模型错误 & 碰撞中误用机械能守恒 & 写出模型条件卡 \\
受力错误 & 漏摩擦、支持力方向错 & 重画受力图 \\
方向错误 & 矢量正负号混乱 & 先规定正方向 \\
计算错误 & 单位、指数、根号错误 & 单独做计算复查 \\
表达错误 & 过程写不清、结论无单位 & 按规范步骤重写 \\
\bottomrule
\end{tabular}
\end{center}

\subsection*{2.\quad 七天复盘安排}

\begin{enumerate}
  \item 第 1 天：复盘运动学、图像和矢量。
  \item 第 2 天：复盘受力分析、摩擦和牛顿定律。
  \item 第 3 天：复盘能量、动量、圆周运动。
  \item 第 4 天：复盘电场、电势、电路。
  \item 第 5 天：复盘磁场、电磁感应和交流电。
  \item 第 6 天：复盘光学、热学、近代物理。
  \item 第 7 天：做一套综合题，只按错因分类，不按分数评价。
\end{enumerate}

\subsection*{3.\quad 综合训练题组}

\begin{enumerate}
  \item \textbf{力学综合：} 斜面、粗糙水平面、弹簧组合，要求写出每一阶段能量方程。
  \begin{center}
  % TikZ 图：粗糙水平面上压缩弹簧
% Day 30：综合训练
\begin{tikzpicture}[scale=0.88,>=Latex,line cap=round,line join=round]
  \path[use as bounding box] (-0.35,-0.82) rectangle (7.80,2.08);
  \draw[thick] (0,0) -- (7.35,0);
  \foreach \x in {0.25,0.75,...,7.0}{
    \draw[gray] (\x,0) -- ++(0.18,-0.18);
  }
  \draw[thick] (7.35,0) -- (7.35,1.35);
  \draw[decorate,decoration={coil,aspect=0.45,segment length=4pt,amplitude=4pt},thick]
    (5.35,0.55) -- (7.35,0.55);

  % 刚接触弹簧时的位置（虚线）与最大压缩位置（实线）分开画。
  \draw[dashed,gray!80,fill=gray!6] (3.35,0.18) rectangle (4.15,0.92);
  \node[gray!80] at (3.75,0.55) {$m$};
  \draw[->,red,thick] (3.35,1.16) -- (4.15,1.16)
    node[midway,above] {$v_0$};

  \draw[fill=gray!14,thick] (4.55,0.18) rectangle (5.35,0.92);
  \node at (4.95,0.55) {$m$};
  \draw[->,orange,thick] (5.25,1.43) -- (4.35,1.43)
    node[midway,above] {$f=\mu mg$};

  \draw[dashed,blue!65] (4.15,0.08) -- (4.15,-0.35);
  \draw[dashed,blue!65] (5.35,0.08) -- (5.35,-0.35);
  \draw[<->,blue,thick] (4.15,-0.42) -- (5.35,-0.42)
    node[midway,below] {$x$};
  \node[font=\scriptsize,align=center] at (2.08,1.55)
    {虚线：刚接触弹簧};
  \node[font=\scriptsize,align=center] at (6.62,1.58)
    {最大压缩\\$v=0$};
\end{tikzpicture}

  \end{center}
  \item \textbf{动量综合：} 碰撞后进入圆轨道，要求区分碰撞阶段和圆周阶段。
  \item \textbf{电场综合：} 带电粒子先加速后偏转，要求画出速度分解图。
  \begin{center}
  % TikZ 图：先加速后偏转模型
% Day 17：带电粒子在电场中的运动
\begin{tikzpicture}[scale=0.78,>=Latex]
  % 加速区
  \draw[thick] (-3.2,1.0) -- (-3.2,-1.0);
  \draw[thick] (-1.6,1.0) -- (-1.6,-1.0);
  \node[above] at (-3.2,1.04) {$+$};
  \node[above] at (-1.6,1.04) {$-$};
  \foreach \y in {-0.45,0,0.45} {
    \draw[->,blue] (-3.0,\y) -- (-1.8,\y);
  }
  \draw[<->] (-3.2,1.32) -- (-1.6,1.32);
  \node[above,blue,fill=white,inner sep=1pt] at (-2.4,1.32) {$U_1$};
  \node[below left,red,fill=white,inner sep=1pt] at (-3.6,0) {$+q$};
  \draw[red,thick,dashed] (-3.6,0) -- (-1.62,0);
  \draw[-{Stealth[length=5pt,width=6pt]},red,line width=1.05pt] (-1.35,0) -- (-0.35,0);
  \node[red,fill=white,inner sep=1pt] at (-0.82,0.28) {$v_0$};

  % 偏转区
  \draw[thick] (-0.2,1.2) -- (3.0,1.2);
  \draw[thick] (-0.2,-1.2) -- (3.0,-1.2);
  \node[above] at (-0.35,1.2) {$+$};
  \node[below] at (-0.35,-1.2) {$-$};
  \foreach \x in {0.35,1.05,1.75,2.45} {
    \draw[->,blue] (\x,1.0) -- (\x,-1.0);
  }
  \draw[<->] (-0.2,1.55) -- (3.0,1.55);
  \node[above,fill=white,inner sep=1pt] at (1.4,1.55) {$L$};
  \draw[<->] (3.35,1.2) -- (3.35,-1.2);
  \node[right,fill=white,inner sep=1pt] at (3.35,0) {$d$};
  \node[blue,fill=white,inner sep=1pt] at (0.15,-1.48) {$U_2$};

  \draw[red,thick,dashed,domain=0.25:3.3,samples=80]
    plot (\x,{-0.10*(\x+0.2)*(\x+0.2)});
  \draw[dashed,gray] (3.3,0) -- (3.3,-1.23);
  \draw[<->] (4.0,0) -- (4.0,-1.23);
  \node[right,fill=white,inner sep=1pt] at (4.0,-0.62) {$y$};

  \node[font=\small,fill=white,inner sep=1pt] at (-2.4,-1.48) {加速区};
  \node[font=\small,fill=white,inner sep=1pt] at (2.15,-1.82) {偏转区};
\end{tikzpicture}

  \end{center}
  \item \textbf{电磁感应综合：} 导体棒在磁场中运动，要求从 $v$ 推到 $\mathcal E,I,F_A$。
  \begin{center}
  % TikZ 图：电磁感应综合题
% Day 30：综合训练
\begin{tikzpicture}[scale=0.9,>=Latex]
  \draw[thick] (-2.5,1.0) -- (2.8,1.0);
  \draw[thick] (-2.5,-1.0) -- (2.8,-1.0);
  \draw[thick] (-2.5,1.0) -- (-2.5,0.32);
  \draw[thick,fill=white] (-2.72,0.32) rectangle (-2.28,-0.32);
  \node[font=\small] at (-2.50,0) {$R$};
  \draw[thick] (-2.5,-0.32) -- (-2.5,-1.0);
  \draw[very thick,red] (-0.3,-1.0) -- (-0.3,1.0);
  \foreach \x in {-1.8,-1.0,-0.2,0.6,1.4,2.2} {
    \foreach \y in {-0.55,0.25} {
      \draw[blue] (\x-0.05,\y-0.05) -- (\x+0.05,\y+0.05);
      \draw[blue] (\x-0.05,\y+0.05) -- (\x+0.05,\y-0.05);
    }
  }
  \draw[->,purple,thick] (-0.3,0.45) -- (1.05,0.45) node[right] {$F$};
  \draw[->,red,thick] (-0.3,0.05) -- (0.85,0.05) node[right] {$v$};
  \draw[->,orange,thick] (-0.3,-0.35) -- (-1.35,-0.35) node[left] {$F_A$};
  \draw[->,blue!70!black,thick] (-0.3,0.62) -- (-0.3,0.93)
    node[right,font=\small] {$I$};
  \node[blue] at (3.15,0.8) {$B$ 向里};
  \node[font=\small,align=center] at (0,-1.65)
    {外力做功 $\rightarrow$ 电能 $\rightarrow$ 焦耳热；终端速度由 $F=F_A$ 决定};
\end{tikzpicture}

  \end{center}
  \item \textbf{实验综合：} 给出一组数据，要求作图、求斜率、解释误差；也可以给出光电门或油膜法情境，要求说明直接测量量和间接量。
  \item \textbf{光学综合：} 给出三棱镜和白光，要求画出红光、紫光的偏折位置；再给出双缝、屏与波长数据，判断相邻亮纹间距及改变波长或缝屏距离后的变化。
  \item \textbf{热学综合：} 给出一段 $p$-$V$ 过程，要求判断气体做功正负、内能变化和吸放热。
  \item \textbf{近代物理综合：} 给出光电效应或能级跃迁数据，要求区分频率、光强、截止电压和能级差的作用。
  \item \textbf{科学发展综合：} 选择一种能源技术（如光伏发电、核电、火力发电或风力发电），用“能量来源、转换、传输、储存、利用、损耗”六个词写出分析链条。
  \item \textbf{数量级估算：} 估算一个家庭一天的照明用电量，写出所选取的主要物理量和近似值。
\end{enumerate}

\textbf{最后提醒：} 学完一本讲义不是终点。后续进步来自持续复述题目、说明每一步的物理理由，并用新题检验模型是否真正会用。

\newpage
\section*{九、易错点}
\begin{mistakebox}
\begin{enumerate}
  \item \textbf{不画图：} 开始列式前先补一张最小图：物体、轨道、场区、电路或状态至少画出一个；方向和阶段要在图上标清。
  \item \textbf{全过程只用一个规律：} 先把全过程切成若干阶段。每到转折点，都重新判断受力、守恒条件和约束关系。
  \item \textbf{方程没有物理解释：} 只列公式不说明来源，容易漏项或错项；每个方程旁写规律名称和适用阶段。
  \item \textbf{忽略临界词：} “刚好”“最大”“最小”“恰不离开”通常对应临界条件；把临界状态单独画出并写出等号条件。
  \item \textbf{答案不检验：} 结尾至少做一项复核：量纲是否正确，方向是否符合图示，极限情况是否合理。
\end{enumerate}
\end{mistakebox}

\section*{十、今日小结}
\begin{enumerate}
  \item 综合题要按画图、选对象、分过程、列方程、复核结果推进。
  \item 复杂题常常是多个基本模型的组合；每段只用该段条件允许的规律，并用中间量把相邻阶段接起来。
  \item 机械能、动量、电磁、光学、热学和近代物理都要检查适用条件，尤其要看守恒条件、理想模型条件和忽略项是否成立。
  \item 能源技术题也要建模：先列能量链，再看守恒、效率、损耗和安全边界。
  \item 数量级估算能帮助判断真实问题中的答案是否合理。
  \item 方程必须服务于物理过程，不能只堆符号。
  \item 后续学习要从知识点记忆转向模型复盘和错因复盘：每道错题至少归因到对象选错、阶段分错、条件漏查、图像误读或计算失误中的一类，并写出下次做题前要检查的一个动作。
\end{enumerate}

\textbf{一分钟复盘：} 回看全书目录，给每一大块画一个四列表：研究对象、典型图、优先规律、检查方式。每一行至少写出一个代表题型；若某块说不出典型图和典型模型，就把它列入下一轮复习，并写明下一次先补图像、公式条件还是典型例题。
